\documentclass[12pt, titlepage]{article}
\usepackage{amsmath}
\usepackage{amsfonts}
\usepackage{amssymb}
\usepackage{fancyhdr}
\usepackage{cancel}
\usepackage{xcolor}
\usepackage[bidi=basic]{babel}
\babelprovide[main,import]{hebrew}
\babelfont[hebrew]{rm}{Hadasim CLM}
\babelfont[hebrew]{sf}{Miriam Mono CLM}
\usepackage{titlesec}
\newcommand\Ccancel[2][black]{\renewcommand\CancelColor{\color{#1}}\cancel{#2}}
\title{חשבון אינפי 2 - 01040281 \\ גליון 4}
\author{יונתן אבידור - 214269565}
\titleformat{\section}{\Large\sffamily\bfseries}{\thesection}{1em}{}
\begin{document}
\maketitle
\pagestyle{fancy}
\fancyhead{}
\fancyfoot[L]{\text{214269565}}
\part*{שאלה 1}
הוכיחו שסדרות הפונקציות הבאות מתכנסות במידה שווה בקטעים הבאים
\section*{סעיף א'}
\begin{gather*}
    \left\{\frac{1}{n} \cos \left(n^2x\right) \right\}^\infty_{n=1}\\
    x \in \mathbb{R}
\end{gather*}
\section*{סעיף ב'}
\begin{gather*}
    \left\{n\ln \left(1+\frac{1}{nx}\right)\right\}^\infty_{n=1}\\
    x\in\left[1,4\right]
\end{gather*}
\part*{פתרון 1}
\section*{סעיף א'}
נתחיל בלהוכיח התכנסות נקודתית\\
יהי $x_0 \in \mathbb{R}$, נתבונן בסדרה $\left\{\frac{1}{n} \cos\left(n x_0\right)\right\}_{n=1}^{\infty}$ \\
$\cos\left(n x_0\right)$ חסומה, בעוד $\frac{1}{n} \xrightarrow{n\to\infty}0$, לכן משיקול "חסומה כפול אפסה" הסדרה מתכנסת ל-$0$\\
נראה כעת כי $\left\{\frac{1}{n} \cos \left(n^2x\right) \right\}^\infty_{n=1}$ שואפת ל$f \equiv 0$ במידה שווה \\
יהי $\varepsilon > 0$ ונבחר $N_{\varepsilon} = \frac{1}{\varepsilon}$, לכל $\mathbb{N} \ni n > N_{\varepsilon}$ ולכל $x \in \mathbb{R}$ מתקיים
\[
\left|\frac{1}{n}\cos\left(n^2 x\right) - \underbrace{f\left(x\right)}_{= 0} \right| \leq \frac{1}{n} \underbrace{\left|\cos\left(n^2x\right)\right|}_{\leq 1}\leq \frac{1}{n} < \varepsilon
\]
לכן הסדרה מתכנסת במ"ש ל-$f\left(x\right) = 0$\\
$\blacksquare$
\section*{סעיף ב'}
נתחיל בלהוכיח התכנסות נקודתית\\
יהי $x_0 \in \left[1,4\right]$, נתבונן בסדרה $\left\{n\ln\left(1+\frac{1}{nx_0}\right)\right\}^\infty_{n=1}$
\begin{gather*}
    n\ln\left(1+\frac{1}{nx_0}\right) = \ln\left(\left(1+\frac{1}{nx_0}\right)^n\right)= \ln\left(\left(1+\frac{1}{nx_0}\right)^\frac{nx_0}{x_0}\right) \\
    =  \ln\left(\left(\left(1+\frac{1}{nx_0}\right)^{nx_0}\right)^{\frac{1}{x_0}}\right) = \frac{1}{x_0} \ln\underbrace{\left(\left(1+\frac{1}{nx_0}\right)^{nx_0}\right)}_{\xrightarrow{n\to\infty}e} \xrightarrow{n\to\infty} \frac{1}{x_0} \ln(e) = \frac{1}{x_0}
\end{gather*}
לכן הסדרה מתכנסת \textbf{נקודתית} ל-$\frac{1}{x}$ \\ נוכיח התכנסות במידה שווה באמצעות האפיון השקול\\
נרצה להראות כי
\[
    \lim_{n\to\infty}    \sup_{x\in \left[1,4\right]} \left|f_n\left(x\right) - f\left(x\right)\right| = 0
\]
נגדיר סדרת פונקציות $g_n(x) = f_n(x) -f(x)$
נגזור את הפונקציות $g_n(x)$ (תהליך הגזירה זהה לכל $n$)
\begin{gather*}
g_n'(x)=    \left(f_n(x) - f(x)\right)' \underset{\text{לינאריות הנגזרת}}{=} f_n'\left(x\right) - f'\left(x\right)\\
    = \left(n\ln\left(1+\frac{1}{nx}\right)\right)' -\left(\frac{1}{x}\right)' = n \cdot \frac{1}{1+\frac{1}{nx}} \cdot \left(-\frac{1}{nx^2}\right) + \frac{1}{x^2} \\
    = n \cdot \frac{\cancel{nx}}{1+nx} \cdot \left(-\frac{1}{\cancel{nx^2}}\right) + \frac{1}{x^2} =  -\frac{n}{x\left(1+nx\right)} + \frac{1}{x^2} \\
    = \frac{-nx+1+nx}{x^2\left(1+nx\right)} = \frac{1}{x^2\left(1+nx\right)} 
\end{gather*}
לכל $n > 0$ ולכל $x \in \left[1,4\right]$ $x^2\left(1+nx\right)>0$, $\frac{1}{x^2\left(1+nx\right)}>0$, כלומר הנגזרת עולה ממש\\
נתבונן בהתנהגות של $g_n(x)$ בקצוות\\
כאשר $x=1$
\begin{gather*}
    \lim_{n\to\infty} \left(f_n\left(1\right) - f\left(1\right)\right) = \lim_{n\to\infty} n\ln\left(1+\frac{1}{n}\right) -\frac{1}{1} = \lim_{n\to\infty} \ln\left(\left(1+\frac{1}{n}\right)^n\right) - 1 \\
    =  \ln(e) - 1 =  1 -1 = 0
\end{gather*}
כאשר $x=4$
\begin{gather*}
    \lim_{n\to\infty} \left(f_n\left(4\right) - f\left(4\right)\right) = \lim_{n\to\infty} n\ln\left(1+\frac{1}{4n}\right) -\frac{1}{4} \\
    \underset{4\text{כפל וחלוקה ב-}}{=} \lim_{n\to\infty} \frac{1}{4} \cdot 4n \ln\left(1+\frac{1}{4n}\right) - \frac{1}{4} \\
    =\lim_{n\to\infty} \frac{1}{4} \ln\left(\left(1+\frac{1}{4n}\right)^{4n}\right) - \frac{1}{4} =  \frac{1}{4}\cdot\ln(e) -\frac{1}{4} = \frac{1}{4}-\frac{1}{4} = 0
\end{gather*}
מכיוון ש-$g_n(x)$ מונוטונית עולה ממש ב-$[1,4]$, מתקיים
\[
g_n(1) \leq g_n(x) \leq g_n(4) 
\]
לכן, עבור הערך המוחלט מתקיים
\[
\forall x \in \left[1,4\right]: \left|g_n(x)\right| \leq \max\left(\left|g_n(1)\right|,\left|g_n(4)\right|\right)
\]
מכיוון שזה נכון לכל $x$, בפרט זה נכון עבור הסופרימום
\[
0 \leq \sup_{x\in[1,4]}  \left|g_n(x)\right| \leq \max\left(\left|g_n(1)\right|,\left|g_n(4)\right|\right) \xrightarrow{n\to\infty} 0
\]
לכן מכלל הסנדוויץ'
\[
    \lim_{n\to\infty}    \sup_{x\in[1,4]} \left|g_n(x)\right| = 0
\]
כלומר
\[
    \lim_{n\to\infty}    \sup_{x\in[1,4]} \left|f_n(x) - f(x)\right| = 0
\]
מכאן, מהאפיון השקול להתכנסות במ"ש. $f_n(x)$ מתכנסת במ"ש לפונקציה $\frac{1}{x}$\\
$\blacksquare$
\newpage
\part*{שאלה 2}
תהא
\[
    f_n\left(x\right) = \sqrt{n} x e^{-nx}, \text{for }x\geq 0
\]
\section*{סעיף א'}
הראו שהסדרה $\left\{f_n(x)\right\}^\infty_{n=1}$ מתכנס נקודתית לכל $x\geq 0$
\section*{סעיף ב'}
הראו שהסדרה מתכנסת במ"ש עבור $x \in \left[0,\infty\right)$
\part*{פתרון 2}
\section*{סעיף א'}
יהי $x_0 \in \left[0,\infty\right)$ נראה כי הסדרה $\left\{\sqrt{n}x_0e^{-nx_0}\right\}^\infty_{n=1}$ מתכנסת\\
נחלק למקרים על פי $x_0$
\subsection*{$x_0 = 0$}
כאשר $x_0$ הסדרה היא $\left\{\sqrt{n}\cdot 0 \cdot e^{-n \cdot 0}\right\}^\infty_{n=1}=\left\{0\right\}^\infty_{n=1}$, כלומר הסדרה הקבועה $0$, לכן היא שואפת נקודתית ל-$f\left(x\right) \equiv 0$
\subsection*{$x_0 > 0$}
\[
    \lim_{n\to\infty} \sqrt{n}x_0 e^{-nx_0} = \lim_{n\to\infty} \frac{\sqrt{n}x_0}{e^{nx_0}}
\]
נשים לב כי במונה יש שורש $n$ כפול קבוע, בעוד שבמכנה יש סדרה מעריכית עם חזקת $n$ כפול קבוע. \\
לכן ממבחן המנה נקבל כי $\lim_{n\to\infty} \frac{\sqrt{n}x_0}{e^{nx_0}} =0$ \\
לכן גם עבור $x_0 > 0$ הסדרה שואפת נקודתית ל-$f\left(x\right) \equiv 0$\\
קיבלנו כי לכל $x_0 \in \left[0,\infty\right)$ $f_n\left(x_0\right) \to 0$ לכן הסדרה מתכנסת נקודתית ל-$f\left(x\right) \equiv 0$\\
$\blacksquare$
\section*{סעיף ב'}
נוכיח כי $\left\{f_n\right\}_{n=1}^\infty$ מתכנסת במ"ש ל-$f\left(x\right) \equiv 0$ בקטע $\left[0,\infty\right)$
נוכיח זאת  באמצעות האפיון השקול להתכנסות במ"ש, כלומר נרצה להראות כי
\[
    \lim_{n\to\infty} \sup_{x\in\left[0,\infty\right)} \left|f_n\left(x\right) - f\left(x\right) \right| = 0
\]
ומכיוון ש-$f\left(x\right)\equiv 0$ זה בעצם
\[
    \lim_{n\to\infty} \sup_{x\in\left[0,\infty\right)} \left|\frac{\sqrt{n}x}{e^{nx}} \right| = 0
\]
מכיוון שבקטע $\left[0,\infty\right)$, $f_n(x)$ תמיד חיובית, $|f_n(x)| = f_n(x)$
נגזור את $f_n(x)$
\[
    f_n'(x) = \frac{\sqrt{n}\cdot e^{nx} - n\cdot e^{nx} \cdot \sqrt{n}x}{e^{2nx}}
\]
נמצא שורשים על מנת למצוא נקודות קיצון
\begin{gather*}
    \frac{\sqrt{n}e^{nx} - x\sqrt{n}ne^{nx}}{e^{2nx}} = 0\\
    \underset{\times e^{2nx}}{\implies} \sqrt{n}e^{nx} - x\sqrt{n}ne^{nx} = 0 \\
    \implies \sqrt{n}e^{nx} \left(1-nx\right) = 0 \\
    \underset{\div \sqrt{n}e^{nx}}{\implies} 1-nx = 0
    \implies x= \frac{1}{n}
\end{gather*}
לפי משפט פרמה, הפונקציה יכולה לשנות את התנהגותה (מבחינת ירידה ועלייה) אך ורק ב-$x=\frac{1}{n}$, נבדוק את ההתנהגות בנקודות קרובות
\begin{gather*}
    f_n\left(0\right) = \frac{\sqrt{n}\cdot 0}{e^0} = 0 \\
    f_n\left(\frac{1}{n}\right) = \frac{\sqrt{n}\frac{1}{n}}{e^{n\frac{1}{n}}} = \frac{\frac{\sqrt{n}}{\sqrt{n}\sqrt{n}}}{e}= \frac{1}{e\sqrt{n}}\\
    f_n(1) = \frac{\sqrt{n}\cdot 1}{e^{n\cdot 1}}= \frac{\sqrt{n}}{e^n}
\end{gather*}
נראה שלכל $n$ מתקיים $\frac{1}{e\sqrt{n}} > \frac{\sqrt{n}}{e^n}$\\
לכל $n$, $n < e^{n-1}$, נשתמש בזה
\[
    n < e^{n-1} \underset{\div e^n}{\implies} \frac{n}{e^n} < \frac{1}{e} \underset{\div \sqrt{n}}{\implies} \frac{\sqrt{n}}{e^n} < \frac{1}{e\sqrt{n}}
\]
לכן הנגזרת עולה ב-$\left[0,\frac{1}{n}\right)$ ויורדת ב-$\left(\frac{1}{n},\infty\right)$\\
לכן הסופרימום נמצא ב-$x=\frac{1}{n}$,
נמצא את הגבול של סדרת הפונקציות בנקודה זו
\[
    \lim_{n\to\infty} \frac{\sqrt{n}\frac{1}{n}}{e^{n\frac{1}{n}}} = \lim_{n\to\infty} \frac{\frac{\sqrt{n}}{\sqrt{n}\sqrt{n}}}{e} = \lim_{n\to\infty} \frac{1}{e\sqrt{n}} = 0
\]
קיבלנו שהסופרימום שואף ל-$0$, לכן $f_n(x)$ מתכנסת במידה שווה ל-$f\left(x\right)$ בקטע $\left[0,\infty\right)$
\newpage
\part*{שאלה 3}
\section*{סעיף א'}
הראו שהסדרה $\left\{\frac{nx}{1+n^2x^2}\right\}_{n=1}^\infty$ מתכנסת במידה שווה בקטע $\left[a,\infty\right)$ לכל $a>0$, אבל מתכנסת נקודתית בלבד ב-$\left(0,\infty\right)$
\section*{סעיף ב'}
הראו שהסדרה $\left\{\frac{1}{n}\ln\left(1+nx\right)\right\}^\infty_{n=1}$ מתכנסת במידה שווה בקטע $\left[0,b\right]$ לכל $b>0$ אבל מתכנסת נקודתית בלבד ב-$\left[0,\infty\right)$
\part*{פתרון 3}
\section*{סעיף א'}
יהי $a>0$
נטען כי הסדרה $\left\{\frac{nx}{1+n^2x^2}\right\}_{n=1}^\infty$ מתכנסת במ"ש לפונקציה $f\left(x\right) \equiv 0$ בקטע $\left[a,\infty\right)$\\
יהי $\varepsilon > 0$, נגדיר $N = \lfloor \frac{1}{\varepsilon a} + 1 \rfloor$\\
לכל $n>N$ ולכל $x\in\left[a,\infty\right)$
\[
    \left|\frac{nx}{1+n^2x^2} - 0\right| = \frac{nx}{1+n^2x^2} < \frac{nx}{n^2x^2} = \frac{1}{nx} < \frac{1}{na} < \frac{1}{Na} = \varepsilon
\]
לכן לכל $a>0$, הסדרה מתכנסת שווה ל-$0$ בקטע $\left[a,\infty\right)$\\
כעת נוכיח כי היא \textbf{לא} מתכנסת במ"ש בקטע $\left(0,\infty\right)$
\section*{סעיף ב'}
\newpage
\part*{שאלה 4}
מצאו את התחום שבו הטור הבא מתכנס בצורה שווה והתחום שבו הוא מתכנס נקודתית
\[
    \sum_{n=1}^\infty \frac{\left(-1\right)^n}{n}e^{n\left(x^2-3x+2\right)}
\]
\part*{פתרון 4}
\newpage
\part*{שאלה 5}
תהי $f_0\left(x\right)$ אינטגרבילית ב-$\left[0,a\right]$. לכל $n\geq 1$ נגדיר $f_n\left(x\right)=\int^x_0 f_{n-1}(t)\,dt, x\in\left[0,a\right]$. הוכיחו כי הסדרה $\left\{f_n(x)\right\}_{n=1}^\infty$ מתכנסת ל-$0$ במידה שווה בקטע $\left[0,a\right]$
\part*{פתרון 5}
\newpage
\part*{שאלה 6}
תהא
\[
    f_n(x) = \frac{nx}{1+n^2x^p}, x\in[0,1], p>0
\]
\section*{סעיף א'}
עבור אילו ערכי $p$ הסדרה $\left\{f_n(x)\right\}^\infty_{n=1}$ מתכנסת במידה שווה ל-$0$
\section*{סעיף ב'}
האם $\lim_{n\to\infty} \int_0^1 f_n(x)\,dx =0$ עבור $p=2$? עבור $p=4$?
\part*{פתרון 6}
\section*{סעיף א'}
\section*{סעיף ב'}
\newpage
\part*{שאלה 7}
מצאו את רדיוסי ההתכנסות של טורי החזקות הבאים
\section*{סעיף א'}
\[
    \sum_{n=1}^\infty \frac{\left(-1\right)}{n!} \left(\frac{n}{e}\right)^n x^n
\]
\section*{סעיף ב'}
\[
    \sum_{n=1}^\infty \frac{1}{n} \left(1+\frac{1}{n}\right)^{n^2} x^{3n}
\]
\part*{פתרון 7}
\section*{סעיף א'}
\section*{סעיף ב'}
\newpage
\part*{שאלה 8}
הוכיחו כי עבור $\left|x\right| < 1$
\[
    \frac{1}{\left(1-x\right)^3} = \frac{1}{2} \sum_{n=0}^{\infty} \left(n+1\right)\left(n+2\right) x^n
\]
\part*{פתרון 8}
\end{document}
